
\documentclass{article}
\usepackage{colortbl}
\usepackage{makecell}
\usepackage{multirow}
\usepackage{supertabular}

\begin{document}

\newcounter{utterance}

\centering \large Interaction Transcript for game `hot\_air\_balloon', experiment `air\_balloon\_survival\_en\_negotiation\_hard', episode 3 with Qwen3{-}32B{-}t0.0.
\vspace{24pt}

{ \footnotesize  \setcounter{utterance}{1}
\setlength{\tabcolsep}{0pt}
\begin{supertabular}{c@{$\;$}|p{.15\linewidth}@{}p{.15\linewidth}p{.15\linewidth}p{.15\linewidth}p{.15\linewidth}p{.15\linewidth}}
    \# & $\;$A & \multicolumn{4}{c}{Game Master} & $\;\:$B\\
    \hline

    \theutterance \stepcounter{utterance}  
    & & \multicolumn{4}{p{0.6\linewidth}}{
        \cellcolor[rgb]{0.9,0.9,0.9}{
            \makecell[{{p{\linewidth}}}]{
                \texttt{\tiny{[P1$\langle$GM]}}
                \texttt{You are participating in a collaborative negotiation game.} \\
\\ 
\texttt{Together with another participant, you must agree on a single set of items that will be kept. Each of you has your own view of how much each item matters to you (importance). You do not know how the other participant values the items. Additionally, you are given the effort each item demands.} \\
\texttt{You may only agree on a set if the total effort of the selected items does not exceed a shared limit:} \\
\\ 
\texttt{LIMIT = 3681} \\
\\ 
\texttt{Here are the individual item effort values:} \\
\\ 
\texttt{Item effort = \{"C17": 700, "B41": 586, "A71": 505, "C00": 116, "B35": 664, "A77": 940, "A03": 387, "A04": 392, "C54": 209, "C01": 571, "A87": 4, "B32": 285, "B86": 651, "B33": 613, "A10": 740\}} \\
\\ 
\texttt{Here is your personal view on the importance of each item:} \\
\\ 
\texttt{Item importance values = \{"C17": 138, "B41": 583, "A71": 868, "C00": 822, "B35": 783, "A77": 65, "A03": 262, "A04": 121, "C54": 508, "C01": 780, "A87": 461, "B32": 484, "B86": 668, "B33": 389, "A10": 808\}} \\
\\ 
\texttt{Goal:} \\
\\ 
\texttt{Your goal is to negotiate a shared set of items that benefits you as much as possible (i.e., maximizes total importance to YOU), while staying within the LIMIT. You are not required to make a PROPOSAL in every message {-} you can simply negotiate as well. All tactics are allowed!} \\
\\ 
\texttt{Interaction Protocol:} \\
\\ 
\texttt{You may only use the following structured formats in a message:} \\
\\ 
\texttt{PROPOSAL: \{'A', 'B', 'C', …\}} \\
\texttt{Propose keeping exactly those items.} \\
\texttt{REFUSE: \{'A', 'B', 'C', …\}} \\
\texttt{Explicitly reject opponent's proposal.} \\
\texttt{ARGUMENT: \{'...'\}} \\
\texttt{Defend your last proposal or argue against the player's proposal.} \\
\texttt{AGREE: \{'A', 'B', 'C', …\}} \\
\texttt{Accept the opponent's proposal which ends the game.} \\
\texttt{STRATEGIC REASONING: \{'...'\}} \\
\texttt{	Describe your strategic reasoning or anticipation explaining your choice of action. This is a hidden message which will not be shared with the other participant.} \\
\\ 
\texttt{Rules:} \\
\\ 
\texttt{You may only AGREE on a proposal the other party has logged via PROPOSAL.} \\
\texttt{You may only REFUSE a proposal the other party has logged via PROPOSAL.} \\
\texttt{Total effort of any PROPOSAL or AGREE set must be ≤ LIMIT.} \\
\texttt{Do NOT reveal your hidden importance scores.} \\
\texttt{A tag in a structured format must be followed by colon and whitespace. The argument must be a python set containing 0 or more strings.} \\
\texttt{So, it must be of the form TAG: \{...\}} \\
\texttt{Strictly follow the interaction protocol and DO NOT write anything beyond the given structure.} \\
\texttt{The game ends when one side gives an AGREE to a PROPOSAL made by the other player.} \\
\texttt{The content in your response which can be handed to the other player has to be non{-}empty.} \\
\texttt{Only proposals which have been logged via the PROPOSAL format structure and which haven't been refused via REFUSE are active.} \\
\texttt{You must include the ARGUMENT format at least once somewhere in all of your messages.} \\
\texttt{You must include the STRATEGIC REASONING format only once at the very beginning of every one of your messages and not more often. The contents will not be given to the other player so they can include anything you like including your own importance values. Here you should reason step by step to come up with you next move.} \\
            }
        }
    }
    & \\ \\

    \theutterance \stepcounter{utterance}  
    & \multicolumn{4}{p{0.6\linewidth}}{
        \cellcolor[rgb]{0.9,0.9,0.9}{
            \makecell[{{p{\linewidth}}}]{
                \texttt{\tiny{[P1$\rangle$GM]}}
                \texttt{<think>} \\
\texttt{Okay, let's see. I need to figure out the best set of items to propose that maximizes my total importance without exceeding the effort limit of 3681. First, I should list out all the items with their effort and my importance values.} \\
\\ 
\texttt{Looking at the data:} \\
\\ 
\texttt{Item Effort and My Importance:} \\
\texttt{{-} C17: 700 (138)} \\
\texttt{{-} B41: 586 (583)} \\
\texttt{{-} A71: 505 (868)} \\
\texttt{{-} C00: 116 (822)} \\
\texttt{{-} B35: 664 (783)} \\
\texttt{{-} A77: 940 (65)} \\
\texttt{{-} A03: 387 (262)} \\
\texttt{{-} A04: 392 (121)} \\
\texttt{{-} C54: 209 (508)} \\
\texttt{{-} C01: 571 (780)} \\
\texttt{{-} A87: 4 (461)} \\
\texttt{{-} B32: 285 (484)} \\
\texttt{{-} B86: 651 (668)} \\
\texttt{{-} B33: 613 (389)} \\
\texttt{{-} A10: 740 (808)} \\
\\ 
\texttt{Hmm, the goal is to maximize my total importance. So I need to prioritize items that give me the highest importance per unit of effort. Let me calculate the ratio for each item.} \\
\\ 
\texttt{Calculating importance per effort:} \\
\\ 
\texttt{{-} A71: 868 / 505 ≈ 1.718} \\
\texttt{{-} C00: 822 / 116 ≈ 7.086} \\
\texttt{{-} A10: 808 / 740 ≈ 1.092} \\
\texttt{{-} C54: 508 / 209 ≈ 2.430} \\
\texttt{{-} A87: 461 / 4 = 115.25 (super high)} \\
\texttt{{-} B41: 583 / 586 ≈ 0.995} \\
\texttt{{-} B35: 783 / 664 ≈ 1.179} \\
\texttt{{-} C17: 138 / 700 ≈ 0.19} \\
            }
        }
    }
    & & \\ \\

    \theutterance \stepcounter{utterance}  
    & & \multicolumn{4}{p{0.6\linewidth}}{
        \cellcolor[rgb]{0.9,0.9,0.9}{
            \makecell[{{p{\linewidth}}}]{
                \texttt{\tiny{[P1$\langle$GM]}}
                \texttt{Your response did not start with the proper strategic reasoning tag at the very beginning of your response. The very first structured format must be of the form STRATEGIC REASONING: \{...\}. Try again.} \\
            }
        }
    }
    & \\ \\

    \theutterance \stepcounter{utterance}  
    & \multicolumn{4}{p{0.6\linewidth}}{
        \cellcolor[rgb]{0.9,0.9,0.9}{
            \makecell[{{p{\linewidth}}}]{
                \texttt{\tiny{[P1$\rangle$GM]}}
                \texttt{<think>} \\
\texttt{Okay, let me start over. The user wants me to make sure the first part is the STRATEGIC REASONING. So first, I need to outline my plan.} \\
\\ 
\texttt{I need to maximize my total importance without exceeding the effort limit. Let me list the items with their importance and effort. Then calculate the importance per effort ratio to prioritize. The highest ratios are A87 (461/4=115.25), C00 (822/116≈7.086), C54 (508/209≈2.43), A71 (868/505≈1.718), B35 (783/664≈1.179), A10 (808/740≈1.092), B41 (583/586≈0.995), etc. The lowest are C17, B33, A77.} \\
\\ 
\texttt{So I should start by including the highest ratio items. A87 is a must because it's so efficient. Then C00, C54, A71, B35, A10, B41. Let's add them up:} \\
\\ 
\texttt{A87: 4} \\
\texttt{C00: 116 → total 120} \\
\texttt{C54: 209 → 329} \\
\texttt{A71: 505 → 834} \\
\texttt{B35: 664 → 1498} \\
\texttt{A10: 740 → 2238} \\
\texttt{B41: 586 → 2824} \\
\texttt{Next is C17: 700 → 3524 (still under 3681)} \\
\texttt{Then B86: 651 → 3524+651=4175 (over)} \\
\texttt{So maybe stop at B41. Total effort 2824. Remaining effort: 3681{-}2824=857. Can we add more?} \\
\\ 
\texttt{Check other high importance items. C17 is next but adding it would take us to 3524. Then maybe B32 (285) → 3524+28} \\
            }
        }
    }
    & & \\ \\

    \theutterance \stepcounter{utterance}  
    & & \multicolumn{4}{p{0.6\linewidth}}{
        \cellcolor[rgb]{0.9,0.9,0.9}{
            \makecell[{{p{\linewidth}}}]{
                \texttt{\tiny{[P1$\langle$GM]}}
                \texttt{Your response did not start with the proper strategic reasoning tag at the very beginning of your response. The very first structured format must be of the form STRATEGIC REASONING: \{...\}. Try again.} \\
            }
        }
    }
    & \\ \\

    \theutterance \stepcounter{utterance}  
    & \multicolumn{4}{p{0.6\linewidth}}{
        \cellcolor[rgb]{0.9,0.9,0.9}{
            \makecell[{{p{\linewidth}}}]{
                \texttt{\tiny{[P1$\rangle$GM]}}
                \texttt{<think>} \\
\texttt{Okay, let me make sure I start with the STRATEGIC REASONING tag first. I need to outline my plan step by step.} \\
\\ 
\texttt{First, I'll list all items with their effort and my importance values. Then calculate the importance per unit effort to prioritize which items to include. The highest ratios are A87 (461/4=115.25), C00 (822/116≈7.086), C54 (508/209≈2.43), A71 (868/505≈1.718), B35 (783/664≈1.179), A10 (808/740≈1.092), B41 (583/586≈0.995), etc. The lowest are C17, B33, A77.} \\
\\ 
\texttt{I should start by including the highest ratio items. A87 is a must because it's so efficient. Then C00, C54, A71, B35, A10, B41. Let's add their efforts:} \\
\\ 
\texttt{A87: 4} \\
\texttt{C00: 116 → total 120} \\
\texttt{C54: 209 → 329} \\
\texttt{A71: 505 → 834} \\
\texttt{B35: 664 → 1498} \\
\texttt{A10: 740 → 2238} \\
\texttt{B41: 586 → 2824} \\
\texttt{Next is C17: 700 → 3524 (still under 3681)} \\
\texttt{Then B86: 651 → 3524+651=4175 (over)} \\
\texttt{So maybe stop at B41. Total effort 2824. Remaining effort: 3681{-}2824=857. Can we add more?} \\
\\ 
\texttt{Check other high importance items. C17 is next but adding it would take us to 3524. Then maybe B32 (285) → 3524+285=3809 (still under). Then B3} \\
            }
        }
    }
    & & \\ \\

    \theutterance \stepcounter{utterance}  
    & & & \multicolumn{2}{p{0.3\linewidth}}{
        \cellcolor[rgb]{0.9,0.9,0.9}{
            \makecell[{{p{\linewidth}}}]{
                \texttt{\tiny{[GM$|$GM]}}
                \texttt{abort} \\
            }
        }
    }
    & & \\ \\

    \theutterance \stepcounter{utterance}  
    & & & \multicolumn{2}{p{0.3\linewidth}}{
        \cellcolor[rgb]{0.9,0.9,0.9}{
            \makecell[{{p{\linewidth}}}]{
                \texttt{\tiny{[GM$|$GM]}}
                \texttt{end game} \\
            }
        }
    }
    & & \\ \\

\end{supertabular}
}

\end{document}
